%...UNIT 1: CONSTANT VELOCITY

\newcommand{\xygraphI}[3]{
\begin{tikzpicture}
    \begin{axis}[height=7.7cm,width=8.5cm,
        axis lines=center,
        ylabel={$y$},
        xlabel={$x$},
        ymin=-10,ymax=10,
        xmin=-10,xmax=10,
        ytick={-10,-8,...,10},
        xtick={-10,-8,...,10},
        grid=both, 
        minor y tick num=1,
        minor x tick num=1,
    ]
        \addplot[thick,domain=-10:10] {#1*x + #2};
        \ifprintanswers
            \addplot[red,only marks,mark=*] coordinates{#3};
        \fi
    \end{axis}
\end{tikzpicture}
}


\newcommand{\distancegraphI}[1]{
    \begin{tikzpicture}
        \begin{axis}[height=6.5cm,width=6.8cm,
            axis lines=left,
            xlabel = {Time (hours)},
            ylabel = {Distance (miles)},
            ymin=0, ymax=100,
            xmin=0, xmax=5,
            ytick={0,10,...,100},
            xtick={0,1,...,5},
            minor x tick num=1,
            grid=both,
            ]
            \addplot[very thick,domain=0:5] {#1*x};
        \end{axis}
    \end{tikzpicture}
}

\newcommand{\positiongraphI}[3]{
    \begin{tikzpicture}
        \begin{axis}[height=5.6cm,width=6.5cm,
            axis lines=left,
            xlabel = {Time (s)},
            ylabel = {Position (m)},
            ymin=0, ymax=20,
            xmin=0, xmax=14,
            ytick={0,4,...,20},
            xtick={0,2,...,14},
            minor y tick num=3,
            minor x tick num=1,
            grid=both,
            ]
            \addplot[very thick,domain=0:14] {#1*x + #2};
            \ifprintanswers
                \addplot[red,only marks,mark=*] coordinates{#3};
            \fi
        \end{axis}
    \end{tikzpicture}

%...AI-CHATBOT PROMPT:
% In an xy plane, the y limits are 0 to 20, and the x limits are 0 to 14. Give me 5 excellent linear equations (y = mx + b) for students to practice slope calculations. Output the results in the form \positiongraphI{m}{b}{coordinates} where coordinates should be in the form (x1,y1)(x2,y2), representing 2 points on the line for calculating slope. 
}

\newcommand{\positionfigureI}[1]{
\begin{center}
    \begin{tikzpicture}[y=1cm,x=1cm]
        \foreach \x in {-5,-4,...,5} {
            \draw[ultra thin] (\x,-1mm) -- (\x,1mm) node[below=2mm] {$\x$};
        }
        \draw[->] (-5,0) -- (5,0) node[below=6mm,pos=0.5] {Position (m)};
        \node[above] at (#1,0) {\Strichmaxerl[2]};
        \ifprintanswers
            \draw[very thick,->,red] (0,0.7) -- ++(#1,0) node[above,pos=0.5] {$\vec{x}$};
        \fi
    \end{tikzpicture}  
\end{center}
}

\newcommand{\positionfigureII}[2]{
\begin{center}
\begin{tikzpicture}[y=1cm,x=1cm]
    \foreach \x in {-5,-4,...,5} {
        \draw[ultra thin] (\x,-1mm) -- (\x,1mm) node[below=3mm] {$\x$};
    }
    \draw[->] (-5,0) -- (5,0) node[below=7mm,pos=0.5] {Position (m)};
    \fill (#1,0) circle (1mm) node[above=1mm] {A};
    \fill (#2,0) circle (1mm) node[above=1mm] {B};
\end{tikzpicture}
\end{center}
}


\newcommand{\positionfigureIII}[4]{
\begin{center}
\begin{tikzpicture}[y=3mm,x=8mm]
    \draw[black!10,xstep=1,ystep=2] (0,0) grid (12,4);
    \foreach \x in {0,1,...,12} {
        \draw[ultra thin] (\x,-1mm) -- (\x,1mm) node[below=3mm] {$\x$};
    }
    \draw[->] (0,0) -- (12,0) node[right=1mm] {$x$ (m)};
    \draw[ultra thick,->,rounded corners=0.7mm] (#1,1)
        -- (#2,1) 
        -- (#2,2) 
        -- (#3,2) 
        -- (#3,3) 
        -- (#4,3);
    \node[fill=black!10] at (#1,-3) {$x_i$};
    \node[above=1mm,fill=black!10] at (#4,4) {$x_f$};
\end{tikzpicture}
\end{center}
}

\newcommand{\positionfigureIV}[2]{
    \begin{center}
    \begin{tikzpicture}[y=1.4mm,x=1.4mm]
        \foreach \x in {-50,-40,...,50} {
            \draw[ultra thin] (\x,-1mm) -- (\x,1mm) node[below=3mm] {$\x$};
        }
        \foreach \x in {-45,-35,...,45} {
            \draw[ultra thin] (\x,-0.8mm) -- (\x,0.8mm);
        }
        \draw[->] (-50,0) -- (50,0) node[below=7mm,pos=0.5] {Position (m)};
        \fill (#1,0) circle (1mm) node[above=1mm] {A};
        \fill (#2,0) circle (1mm) node[above=1mm] {B};
    \end{tikzpicture}
    \end{center}
}

\newcommand{\positionfigureV}[6]{
    %...PRO TIP: Type the following into ChatGPT:
    % Generate 5 integers between -10 to 10 (inclusive) such that the absolute value of the sum of the 5 integers is less than or equal to 10. Print each of the integers in each of the first 5 curly brackets below. Calculate the maximum displacement from equilibrium and, if it's greater than 10, enter it in the 6th curly bracket; otherwise enter 10 for the 6th curly bracket. Do not output in markdown; I just want the verbatim text: \positionfigureV{}{}{}{}{}{}
    
    \begin{tikzpicture}[x=4mm,y=4mm]
        \draw[step=1,lightgray] (-10,0) grid (10,4);
        \draw[->,line width=0.6mm,-{Triangle}] (0,4) -- ++(#5,0);
        \draw[->,line width=0.6mm,-{Triangle}] (0,3) -- ++(#4,0);
        \draw[->,line width=0.6mm,-{Triangle}] (0,2) -- ++(#3,0);
        \draw[->,line width=0.6mm,-{Triangle}] (0,1) -- ++(#2,0);
        \draw[->,line width=0.6mm,-{Triangle}] (0,0) -- ++(#1,0);
        \node[below] at (current bounding box.south) {Figure 1: Displacement vectors.};
    \end{tikzpicture}
    
    \smallskip
    
    \begin{tikzpicture}[x=4mm,y=4mm]
        \draw[step=1,lightgray] (-#6,0) grid (#6,4);
        \draw (0,0) node {\scriptsize $|$} node[below=2pt] {$0$};
            \begin{scope}[red,opacity={\ifprintanswers 1\else 0\fi}]
                \draw[->,line width=0.3mm,-{Triangle}] ({#1+#2+#3+#4},4) -- ++(#5,0);
                \draw[->,line width=0.3mm,-{Triangle}] ({#1+#2+#3},3) -- ++(#4,0);
                \draw[->,line width=0.3mm,-{Triangle}] ({#1+#2},2) -- ++(#3,0);
                \draw[->,line width=0.3mm,-{Triangle}] (#1,1) -- ++(#2,0);
                \draw[->,line width=0.3mm,-{Triangle}] (0,0) -- ++(#1,0);
                \pgfmathsetmacro{\vecsum}{#1+#2+#3+#4+#5}
        
                \pgfmathtruncatemacro{\iszero}{(\vecsum==0)}
                
                \ifnum\iszero=1
                    \node[above] at (0,4.6) {vector sum is zero};
                \else
                    \draw[->,line width=0.6mm,-{Triangle}] (0,4.8) -- ++(\vecsum,0)
                        node[above=1mm,pos=0.5] {vector sum};
                \fi
        \end{scope}
        \node[below] at (current bounding box.south) {Figure 2: Head-to-tail addition and vector sum.};
    \end{tikzpicture}
}

\newcommand{\positionfigureVI}[1]{
\begin{tikzpicture}[y=6mm,x=6mm]
    \draw[black!20,step=1] (0,0) grid (24,2);  
    \draw[|-|] (0,-0.3) -- ++(1,0) node[below] {\small 1 meter};
    \ifprintanswers
    \foreach \x in {0,#1,...,24}{
        \fill[red] (\x,1) circle (0.9mm);
    }
    \fi
\end{tikzpicture}
}

\newcommand{\positiongraphIII}[6]{
    %...Generates a quadrant I position vs. time with linear function.
    % #1 = velocity (slope)
    % #2 = initial position (y-intercept)
\begin{tikzpicture}
\begin{axis}[height=5.9cm,width=6.4cm,
    axis lines=left,
    xlabel = {Time (s)},
    ylabel = {Position (m/s)},
    ymin=0, ymax=10,
    xmin=0, xmax=10,
    ytick={0,1,...,10},
    xtick={0,1,...,10},
    grid = both,
    clip=true,
    legend style={at={(1.1,0.85)}}
    % set layers
    ]
    \addplot[very thick,black,domain=0:10] {#1*\x + #2};
    \addlegendentry{A} % <-- Legend for the first plot
    
    \addplot[densely dashed,very thick,black,domain=0:10] {#3*\x + #4};
    \addlegendentry{B} 

    \addplot[only marks, mark=*,mark size=1.5pt,black,domain=0:10] {#5*\x + #6};
    \addlegendentry{C} 
    \end{axis}
\end{tikzpicture}
}

\newcommand{\velocitygraph}[2]{
    \begin{tikzpicture}
        \begin{axis}[height=6cm,width=7cm,
            axis y line=left,
            axis x line=center,
            xlabel = {Time (s)},
            ylabel = {Velocity (m/s)},
            ymin=-5, ymax=5,
            xmin=0, xmax=10,
            ytick={-5,-4,...,5},
            xtick={0,2,...,10},
            grid=both,
            minor x tick num = 1,
            x label style={at={(axis description cs:1,0.5)},right=1mm},
            clip=false
        ]
            \addplot[very thick,black,domain=0:#2] {#1};
            \ifprintanswers
                \begin{pgfonlayer}{background}
                    \fill[red!12] (0,0) rectangle (#2,#1);
                \end{pgfonlayer}
            \fi
            \end{axis}
    \end{tikzpicture}
}

%...AI-CHATBOT PROMPT:
%Generate a random integer between -5 to 5 (inclusive), and another one between 5 and 10 (inclusive). Output the results in the form \velocitygraph{a}{b} where a and b are the generated numbers.




%..UNIT 2: ACCELERATION

\newcommand{\velocitytableI}[4]{%
    % #1 = initial velocity
    % #2 = acceleration
    % #3 = initial time
    % #4 = change in time
\begin{tabular}{|c|c|}
    \hline
    \thead{\bfseries Velocity (m/s)} & \thead{\bfseries Time (s)} \\ \hline
    $\fpeval{#1 + #2*(#3 + 0*#4)}$ & \fpeval{#3 + 0*#4} \\ \hline
    $\fpeval{#1 + #2*(#3 + 1*#4)}$ & \fpeval{#3 + 1*#4} \\ \hline
    $\fpeval{#1 + #2*(#3 + 2*#4)}$ & \fpeval{#3 + 2*#4} \\ \hline
    $\fpeval{#1 + #2*(#3 + 3*#4)}$ & \fpeval{#3 + 3*#4} \\ \hline
    $\fpeval{#1 + #2*(#3 + 4*#4)}$ & \fpeval{#3 + 4*#4} \\ \hline
\end{tabular}
}



\newcommand{\positiontableI}[5]{%
    % #1 = initial position
    % #2 = initial velocity
    % #3 = initial time
    % #4 = change in time
    % #5 = acceleration
\begin{tabular}{|c|c|}
    \hline
    \thead{\bfseries Position (m)} & \thead{\bfseries Time (s)} \\ \hline
    $\fpeval{#1 + #2*(#3 + 0*#4) + 0.5*#5*(#3 + 0*#4)^2}$ & \fpeval{#3 + 0*#4} \\ \hline
    $\fpeval{#1 + #2*(#3 + 1*#4) + 0.5*#5*(#3 + 1*#4)^2}$ & \fpeval{#3 + 1*#4} \\ \hline
    $\fpeval{#1 + #2*(#3 + 2*#4) + 0.5*#5*(#3 + 2*#4)^2}$ & \fpeval{#3 + 2*#4} \\ \hline
    $\fpeval{#1 + #2*(#3 + 3*#4) + 0.5*#5*(#3 + 3*#4)^2}$ & \fpeval{#3 + 3*#4} \\ \hline
    $\fpeval{#1 + #2*(#3 + 4*#4) + 0.5*#5*(#3 + 4*#4)^2}$ & \fpeval{#3 + 4*#4} \\ \hline
\end{tabular}
}

\NewDocumentCommand{\velocitygraphI}{m m o}{
    % Produces a line on a velocity vs time graph.
    % #1 = slope (acceleration) 
    % #2 = y-intercept (initial velocity)
    % #3 [optional] = sample coordinates (x1,y1)(x2,y2) for finding slope
\begin{tikzpicture}
\begin{axis}[height=6cm,width=6.4cm,
    axis y line=left,
    axis x line=center,
    ylabel = {Velocity (m/s)},
    xlabel = {Time (s)},
    ymin=-8, ymax=8,
    xmin=0, xmax=10,
    ytick={-8,-6,...,8},
    xtick={0,1,...,10},
    grid = both,
    minor y tick num=1,
    ]
    \addplot[very thick,black,domain=0:10] {#1*\x + #2};
    \ifprintanswers
            % \IfValueT checks if the 3rd argument was actually provided
            \IfValueT{#3}{
                \addplot[red,only marks,mark=*] coordinates{#3};
            }
        \fi
    \end{axis}
\end{tikzpicture}

%...AI CHATBOT PROMPT:
% A velocity vs. time graph has y limits -8 to 8 and x limits 0 to 10. Generate a linear function for students to practice calculating acceleration. Output the result in the form \velocitygraphI{m}{b}{(x1,y1)(x2,y2)} where m is the slope (acceleration), b is the y-intercept (initial velocity), and (x1,y1)(x2,y2) are two coordinates on the line used to calculate slope. Negative accelerations are acceptable.
}

\newcommand{\positiongraphII}[3]{
    % Generates a position vs. time graph for uniform acceleration motion.
    %...Inspired by https://ophysics.com/k4.html
    % #1 = initial position
    % #2 = initial velocity
    % #3 = acceleration
\begin{tikzpicture}
    \begin{axis}[height=6cm,width=5cm,
        axis y line=left,
        axis x line=center,
        ylabel={Position (m)},
        xlabel={Time (s)},
        ymin=-4,ymax=4,
        xmin=0,xmax=4,
        ytick={-4,-3,...,4},
        xtick={0,1,...,4},
        grid=both,
        x label style={at={(axis description cs: 1,0.5)},right=1mm},
    ]
        \addplot[very thick,domain=0:4] {#1 + #2*x + 0.5*#3*x^2};        
    \end{axis}
\end{tikzpicture}    
}

\newcommand{\velocitygraphII}[2]{
    %...Generates a minimalistic v vs. t graph.
    % #1 = initial velocity (y-intercept)
    % #2 = acceleration (slope)
    %... See https://ophysics.com/k4.html
\begin{tikzpicture}[y=5mm,x=7mm]
    \draw[->] (0,-4.4) -- (0,4.4) node[left,pos=0.9] {$v$} node[pos=0.5,left] {\footnotesize $0$};
    \draw[->] (0,0) -- (4.4,0) node[right] {$t$};
    \draw[very thick,domain=0:4] plot(\x,{#1 + #2*\x});
\end{tikzpicture}
}


\newcommand{\velocitygraphIII}[3]{
    %...Generates a quadrant I velocity vs. time with linear function.
    % #1 = acceleration (slope)
    % #2 = initial velocity (y-intercept)
    % #3 = final time (where line stops)
\begin{tikzpicture}
\begin{axis}[height=6cm,width=6.4cm,
    axis lines=left,
    xlabel = {Time (s)},
    ylabel = {Velocity (m/s)},
    ymin=0, ymax=10,
    xmin=0, xmax=10,
    ytick={0,1,...,10},
    xtick={0,1,...,10},
    grid = both,
    clip=true,
    % set layers
    ]
    \addplot[very thick,black,domain=0:#3] {#1*\x + #2};
    \ifprintanswers
    \begin{pgfonlayer}{background}
        \fill[red!10] (0,#2) -- (#3,{#1*#3+#2}) -- (#3,0) -- (0,0) -- cycle;
    \end{pgfonlayer}
    \fi
    \end{axis}
\end{tikzpicture}

%...AI CHATBOT PROMPT:
% A velocity vs. time graph has y limits 0 to 10 and x limits 0 to 10. Generate a linear function for students to practice calculating displacement (as area under the curve). Output the result in the form \velocitygraphIII{m}{b} where m is the slope (acceleration) and b is the y-intercept (initial velocity). Negative accelerations are acceptable.
}

\newcommand{\triangleI}[3]{%
    % #1 = base
    % #2 = height
    % #3 = a parameter to control figure size
    \begin{tikzpicture}[y=1mm*#3,x=1mm*#3]
        \draw[thick,fill=black!10] (-0.5*#1,0) -- (0.5*#1,0) -- (0,#2) -- cycle;
        \draw[densely dashed] (0,0) -- (0,#2) node[pos=0.4,right] {#2};
        \draw[|-|] (-0.5*#1,-2mm) -- (0.5*#1,-2mm) node[below,pos=0.5] {#1};
    \end{tikzpicture}
}

\newcommand{\triangleII}[3]{%
    % #1 = base
    % #2 = height
    % #3 = a parameter to control figure size
    \begin{tikzpicture}[y=1mm*#3,x=1mm*#3]
        \draw[thick,fill=black!10] (0,0) -- (#1,0) node[below,pos=0.5] {#1} -- (#1,#2) node[right,pos=0.5] {#2} -- cycle;
    \end{tikzpicture}
}